\begin{theorem}\label{thm:hybrid-endpoint-continuity-counterexample}
Fix a hybrid radius \(r_0\in(0,1)\).  There exist a smooth trivial
polarized degeneration
\[
 \pi\colon X=\mathbf P^1_{\mathbf C}\times\mathbb D^*
       \longrightarrow\mathbb D^*,
 \qquad
 L=\operatorname{pr}_1^*\mathcal O_{\mathbf P^1}(1),
\]
and a smooth semipositive metric \(\phi\) on \(L\), with logarithmic
growth at the origin, having the following properties.
\begin{enumerate}
\item The metric \(\phi\) extends as an everywhere finite psh metric on
the product model
\[
 \mathcal X=\mathbf P^1_{\mathbf C}\times\mathbb D,
 \qquad
 \mathcal L=\operatorname{pr}_1^*\mathcal O_{\mathbf P^1}(1).
\]
\item Its associated non-Archimedean metric is the continuous product
Fubini--Study model metric:
\[
                   \phi^{\NA}
                   =\phi_{\mathrm{FS}}^{\NA}.
\]
\item Relative to the product Fubini--Study metric, \(\phi\) has zero
fibrewise supremum on every complex fibre and on the non-Archimedean
fibre:
\[
 \sup_{X_t}(\phi_t-\phi_{\mathrm{FS},t})=0
 \quad(t\ne0),
 \qquad
 \sup_{X_{\mathbf C((t))}^{\mathrm{an}}}
 (\phi^{\NA}-\phi_{\mathrm{FS}}^{\NA})=0.
\]
\item The canonical hybrid psh metric \(\phi^{\mathrm{hyb}}\), obtained
by putting \(\phi\) on the complex fibres and
\(\phi^{\NA}\) on the non-Archimedean fibre, is not continuous.
More precisely, for \(p_*=[0:1]\in\mathbf P^1(\mathbf C)\), there are
two sequences \(x_m^+\) and \(x_m^-\) in the complex locus that both
converge to \((p_*)_{\mathbf C((t))}\) and for which
\[
\begin{aligned}
 \lim_{m\to\infty}
 \bigl(\phi^{\mathrm{hyb}}
       -\phi_{\mathrm{FS}}^{\mathrm{hyb}}\bigr)(x_m^+)
   &=\log r_0,\\
 \lim_{m\to\infty}
 \bigl(\phi^{\mathrm{hyb}}
       -\phi_{\mathrm{FS}}^{\mathrm{hyb}}\bigr)(x_m^-)
   &=0.
\end{aligned}
\]
\end{enumerate}
Consequently, continuity of a semipositive metric on the punctured total
space together with continuity of its associated non-Archimedean metric
does not imply continuity of the canonical hybrid psh extension, even
under the displayed zero-supremum normalizations.

Moreover, relative to the product Fubini--Study metric, the resulting
potential satisfies separately all three conditions in Favre's
Question~1: it is continuous near every punctured
\(\overline X_r\), its curvature is the restriction of a positive
closed \((1,1)\)-current on a model neighborhood, and its
generic-Lelong function extends continuously to the non-Archimedean
analytification.  Nevertheless, that potential is not uniform in the
sense of \cite[Definition~4.1]{Favre20Hybrid}.
\end{theorem}

\begin{proof}
For \(j\geq1\), put
\[
 s_j=2^j,\qquad
 \rho_j=\exp(-s_j),\qquad
 c_j=4^{-j},\qquad
 \epsilon_j=\exp(-s_j/c_j)=\exp(-8^j),
\]
and define on the unit disk
\[
 q_j(t)=\frac{c_j}{2}
   \log\frac{|t-\rho_j|^2+\epsilon_j^2}{5}.
\tag{66.20d}\label{eq:p18-qj}
\]
The functions \(q_j\) are smooth and subharmonic.  Indeed,
\(\log(|t-a|^2+\epsilon^2)\) has nonnegative
\(\mathrm d\mathrm d^c\), and multiplication by \(c_j/2>0\)
preserves this property.  Moreover \(q_j\leq0\) on \(\mathbb D\):
for \(|t|<1\), the numerator in \eqref{eq:p18-qj} is strictly smaller
than \(5\).

Set
\[
                         u(t)=\sum_{j=1}^{\infty}q_j(t).
\tag{66.20e}\label{eq:p18-u}
\]
Let \(K\Subset\mathbb D^*\).  For all sufficiently large \(j\), the
points \(\rho_j\) and the disks of radius \(\epsilon_j\) about them
have a fixed positive distance from \(K\).  Every derivative, on \(K\),
of the logarithm occurring in \eqref{eq:p18-qj} is then bounded by a
constant depending only on \(K\) and the order of the derivative.
Since \(\sum_jc_j<\infty\), the series in \eqref{eq:p18-u} converges in
\(C^\infty(K)\).  Hence \(u\) is smooth and subharmonic on
\(\mathbb D^*\).

At the origin one has
\[
\begin{aligned}
 q_j(0)
 &=
 \frac{c_j}{2}
 \left(
  2\log\rho_j-\log5+
  \log\left(1+\frac{\epsilon_j^2}{\rho_j^2}\right)
 \right)\\
 &=-c_js_j-\frac{c_j}{2}\log5+O(c_j)
  =-2^{-j}+O(4^{-j}).
\end{aligned}
\tag{66.20f}\label{eq:p18-at-zero}
\]
Thus \(\sum_jq_j(0)>-\infty\).  The partial sums in
\eqref{eq:p18-u} form a decreasing sequence of subharmonic functions on
\(\mathbb D\), because every \(q_j\leq0\).  Their limit is not
identically \(-\infty\), by \eqref{eq:p18-at-zero}; therefore the
decreasing-limit theorem gives a subharmonic extension of \(u\) to
\(\mathbb D\), with the finite value prescribed by
\eqref{eq:p18-at-zero}.

This extension has zero Lelong number at the origin.  To see this without
using any pointwise radial assertion, let \(M_u(r)\) be its circular
mean.  The submean inequality and \(u\leq0\) give
\[
                         u(0)\leq M_u(r)\leq0.
\]
The circular-mean characterization of the one-variable Lelong number
therefore yields
\[
                       \nu_0(u)
                       =\lim_{r\to0}\frac{M_u(r)}{\log r}=0.
\tag{66.20g}\label{eq:p18-zero-lelong}
\]

We next establish the logarithmic-growth estimate and isolate two
different normalized limits.  Let
\[
                B_j=\{\,t\in\mathbb D:
                         |t-\rho_j|<\rho_j/2\,\}.
\]
The disks \(B_j\) are pairwise disjoint.  Indeed,
\(\rho_{j+1}=\rho_j^2\), \(\rho_1<1/3\), and hence
\[
                    \frac32\rho_{j+1}
                    <\frac12\rho_j .
\]
If \(t\notin B_j\), then
\[
\begin{aligned}
 q_j(t)
 &\geq\frac{c_j}{2}
       \log\frac{\rho_j^2}{20}\\
 &=-c_js_j-\frac{c_j}{2}\log20.
\end{aligned}
\tag{66.20h}\label{eq:p18-outside-well}
\]
If \(t\) belongs to none of the \(B_j\), summing
\eqref{eq:p18-outside-well} gives \(u(t)\geq-C\), since
\[
                \sum_{j\geq1}c_js_j=\sum_{j\geq1}2^{-j}<\infty.
\]
If \(t\in B_m\), pairwise disjointness permits the use of
\eqref{eq:p18-outside-well} for all \(j\ne m\), while
\[
\begin{aligned}
 q_m(t)
 &\geq\frac{c_m}{2}\log\frac{\epsilon_m^2}{5}\\
 &=-s_m-\frac{c_m}{2}\log5.
\end{aligned}
\tag{66.20i}\label{eq:p18-inside-well}
\]
Moreover \(-\log|t|=s_m+O(1)\) uniformly for \(t\in B_m\).
Equations \eqref{eq:p18-outside-well}--\eqref{eq:p18-inside-well},
together with \(u\leq0\), prove that, after increasing \(C\),
\[
             -\log|t|^{-1}-C\leq u(t)\leq0
             \qquad(0<|t|<1).
\tag{66.20j}\label{eq:p18-log-growth}
\]

At \(t=\rho_m\), the \(m\)-th summand is exactly
\[
                  q_m(\rho_m)
                  =-s_m-\frac{c_m}{2}\log5.
\]
Since \(\rho_m\notin B_j\) for \(j\ne m\), the sum of all other
summands lies between \(-C\) and \(0\), by
\eqref{eq:p18-outside-well}.  It follows that
\[
                         u(\rho_m)=-s_m+O(1).
\tag{66.20k}\label{eq:p18-plus-sequence}
\]
On the other hand, \(-\rho_m\) belongs to none of the \(B_j\), because
\(|-\rho_m-\rho_j|=\rho_m+\rho_j>\rho_j/2\).  Hence
\[
                         -C\leq u(-\rho_m)\leq0.
\tag{66.20l}\label{eq:p18-minus-sequence}
\]
Since \(\log|\rho_m|=-s_m\), equations
\eqref{eq:p18-plus-sequence}--\eqref{eq:p18-minus-sequence} give
\[
\begin{aligned}
 \frac{\log r_0}{\log|\rho_m|}\,u(\rho_m)
   &\longrightarrow\log r_0,\\
 \frac{\log r_0}{\log|-\rho_m|}\,u(-\rho_m)
   &\longrightarrow0.
\end{aligned}
\tag{66.20m}\label{eq:p18-two-limits}
\]

We now put this scalar example on a polarized family without adding it as
a scalar pulled back from the base.  Let
\[
 \mathcal X=\mathbf P^1_{\mathbf C}\times\mathbb D,\qquad
 \mathcal L=\operatorname{pr}_1^*\mathcal O_{\mathbf P^1}(1),
\]
and use homogeneous coordinates \([Z_0:Z_1]\) on \(\mathbf P^1\).
The product Fubini--Study model metric and our new metric are represented
by the homogeneous additive weights
\[
\begin{aligned}
 \phi_{\mathrm{FS}}([Z_0:Z_1],t)
   &=\frac12\log\bigl(|Z_0|^2+|Z_1|^2\bigr),\\
 \phi([Z_0:Z_1],t)
   &=\frac12\log\bigl(|Z_0|^2+\exp(2u(t))|Z_1|^2\bigr).
\end{aligned}
\tag{66.20n}\label{eq:p18n-metrics}
\]
The second formula is smooth on
\(X=\mathcal X|_{\mathbb D^*}\).  It also defines an everywhere finite
psh metric on \(\mathcal L\) over \(\mathcal X\).  Indeed, on the chart
\(Z_1\ne0\), with \(z=Z_0/Z_1\), its local weight is
\[
 F_1(z,t)=\frac12\log\bigl(|z|^2+\exp(2u(t))\bigr),
\tag{66.20o}\label{eq:p18n-weight-one}
\]
which is the logarithmic sum of the exponentials of the two psh
functions \(\log|z|\) and \(u(t)\).  On the chart \(Z_0\ne0\), with
\(w=Z_1/Z_0\), its local weight is
\[
 F_0(w,t)=\frac12\log\bigl(1+\exp(2u(t))|w|^2\bigr),
\tag{66.20p}\label{eq:p18n-weight-zero}
\]
the same construction applied to \(0\) and \(u(t)+\log|w|\).
The convex coordinatewise nondecreasing log-sum-exp function preserves
plurisubharmonicity.  Thus both local weights are psh; their transition
law is the one for \(\mathcal O_{\mathbf P^1}(1)\).  On the punctured
total space they are smooth, so \(\phi\) is a smooth semipositive metric.

Its potential relative to the Fubini--Study metric is the globally
defined function
\[
 \psi([Z_0:Z_1],t)
 =\phi-\phi_{\mathrm{FS}}
 =\frac12\log
   \frac{|Z_0|^2+\exp(2u(t))|Z_1|^2}
        {|Z_0|^2+|Z_1|^2}.
\tag{66.20q}\label{eq:p18n-relative-potential}
\]
Since \(u\leq0\), one has
\[
                         u(t)\leq\psi\leq0.
\tag{66.20r}\label{eq:p18n-potential-bounds}
\]
The upper bound is immediate.  For the lower bound, use
\(|Z_0|^2\geq\exp(2u(t))|Z_0|^2\).  In particular, the extension just
constructed proves logarithmic growth by the first equivalent condition
in \cite[Definition~3.9]{PS23Hybrid}, while
\eqref{eq:p18-log-growth} and \eqref{eq:p18n-potential-bounds} verify
the two-sided logarithmic bound directly.  Moreover,
\(\psi([1:0],t)=0\), and hence
\[
                   \sup_{X_t}\psi_t=0\qquad(t\ne0).
\tag{66.20s}\label{eq:p18n-complex-normalization}
\]

We compute the associated non-Archimedean metric.  Let \(E\) be a prime
vertical divisor on an arbitrary model dominating \(\mathcal X\).  At a
general smooth point of \(E\), one may write
\[
                          t=\zeta^{b_E}h,
\]
where \(\zeta=0\) is a local equation for \(E\), \(h\) is a unit, and
\(b_E=\operatorname{ord}_E(t)\).  In one complex variable the Lelong
number is the atomic mass at the origin of the Riesz measure, up to the
fixed normalization.  Pullback by a holomorphic map of order \(b_E\)
multiplies this atomic mass by \(b_E\).  Thus
\[
                 \nu_E(u\circ\pi)
                 =b_E\nu_0(u)=0
\tag{66.20t}\label{eq:p18n-base-divisorial-zero}
\]
by \eqref{eq:p18-zero-lelong}.

Choose one of the two affine charts containing the image of the general
point of \(E\).  On \(Z_0\ne0\), the psh weight
\eqref{eq:p18n-weight-zero} is bounded below by \(0\), so its generic
Lelong number along \(E\) is zero.  On \(Z_1\ne0\), the psh weight
\eqref{eq:p18n-weight-one} is bounded below by \(u\circ\pi\).
Monotonicity of Lelong numbers, together with
\eqref{eq:p18n-base-divisorial-zero}, again gives
\[
                         \nu_E(\phi)=0.
\tag{66.20u}\label{eq:p18n-metric-divisorial-zero}
\]
The smooth Fubini--Study reference has zero generic Lelong numbers as
well.  \cite[Theorem~3.10]{PS23Hybrid} therefore gives
\[
                  \phi^{\NA}
                  =\phi_{\mathrm{FS}}^{\NA},
\tag{66.20v}\label{eq:p18n-na-model}
\]
the continuous product model metric.  In particular,
\[
 \sup_{X_{\mathbf C((t))}^{\mathrm{an}}}
 (\phi^{\NA}-\phi_{\mathrm{FS}}^{\NA})=0.
\tag{66.20w}\label{eq:p18n-na-normalization}
\]

\cite[Theorem~3.13]{PS23Hybrid} now applies to \(\phi\): the metric and
the endpoint \eqref{eq:p18n-na-model} form one canonical hybrid psh
metric \(\phi^{\mathrm{hyb}}\).  This invokes genuine hybrid psh
gluing, rather than an arbitrary assignment of values on two disjoint
sets.

Put \(p_*=[0:1]\), and write \((p_*)_{\mathbf C((t))}\) for the point
on the non-Archimedean fibre induced by the constant algebraic section.
Functoriality of hybrid analytification, applied to that section, gives
\[
\begin{aligned}
 x_m^+&=(p_*,\rho_m)\longrightarrow(p_*)_{\mathbf C((t))},\\
 x_m^-&=(p_*,-\rho_m)\longrightarrow(p_*)_{\mathbf C((t))}.
\end{aligned}
\tag{66.20x}\label{eq:p18n-hybrid-convergence}
\]
The hybrid metric induced by the product Fubini--Study model is
continuous.  On \(p_*\) the relative potential in
\eqref{eq:p18n-relative-potential} is exactly
\[
                         \psi(p_*,t)=u(t).
\tag{66.20y}\label{eq:p18n-special-section}
\]
In additive weights, the difference on a complex fibre over \(t\ne0\)
is the usual hybrid rescaling
\[
 \bigl(\phi^{\mathrm{hyb}}
       -\phi_{\mathrm{FS}}^{\mathrm{hyb}}\bigr)|_{X_t}
 =\frac{\log r_0}{\log|t|}\,\psi_t,
\tag{66.20z}\label{eq:p18n-hybrid-normalization}
\]
whereas its difference on the central non-Archimedean fibre is zero by
\eqref{eq:p18n-na-model}.  Combining
\eqref{eq:p18-two-limits}, \eqref{eq:p18n-hybrid-convergence},
\eqref{eq:p18n-special-section}, and
\eqref{eq:p18n-hybrid-normalization} proves the two limits in the
statement.  Since \(\log r_0\ne0\), the canonical hybrid psh metric is
not continuous despite the normalizations
\eqref{eq:p18n-complex-normalization} and
\eqref{eq:p18n-na-normalization}.

It remains to check Favre's formulation independently.  Take
\(\varphi=\psi\).  The curvature current
\[
                         T=c_1(\mathcal L,\phi)
\tag{66.20aa}\label{eq:p18n-favre-current}
\]
is positive and closed on the product model because \(\phi\) is a psh
metric, and on the punctured total space it restricts to
\(\omega_{\mathrm{FS}}+\mathrm d\mathrm d^c\varphi\).  The potential
is smooth on \(X\), hence continuous near every punctured
\(\overline X_r\).  Equation \eqref{eq:p18n-metric-divisorial-zero}
says that Favre's generic-Lelong function is \(g_T\equiv0\), so it
extends continuously to the non-Archimedean analytification.  All three
bullets of Favre's Question~1 hold.

Suppose, for contradiction, that \(\varphi\) were uniform.  The remark
following \cite[Definition~4.1]{Favre20Hybrid} states that its normalized
function on the punctured hybrid locus would then extend continuously to
the hybrid space.  For the hybrid radius appearing in that definition,
the two sequences in \eqref{eq:p18n-hybrid-convergence} have the distinct
limits computed above.  Such a continuous extension is impossible.  Thus
\(\varphi\) is not uniform, completing the proof.
\end{proof}
